\begin{abstract}
Predicting how long a lithium-ion cell will last, using only its first few
charge--discharge cycles, would shorten battery development and validation
from months to days. This project follows the data-driven approach of
Severson \textit{et al.}~\cite{severson2019}, using the public dataset of
commercial LFP/graphite cells (A123 APR18650M1A, 1.1\,Ah) cycled to failure
under two-step fast-charging policies. This interim report documents the
experiment and the data, using as the example a batch of 47 cells
(2018-02-20) that is not part of the original study and whose policies all
charge from 0 to 80\% state of charge in 10 minutes.
It describes the conversion of the raw MATLAB files into a Python data
structure and presents an exploratory analysis of capacity fade, internal
resistance, temperature and the discharge voltage curve. Of the 47 cells,
39 have a usable cycle life (392 to 1186 cycles). They fall into two
distinct groups: standard cells (392--514 cycles) and cells marked
\texttt{-newstructure} (777--1186 cycles). The key early-life feature of
\cite{severson2019}, the variance of $\dq$, is strongly correlated with log
cycle life across all cells (Pearson $r=-0.92$). Within each group,
however, the correlation is much weaker ($r=-0.38$ and $r=-0.27$). The
report then describes a data pipeline for all four public batches (163
labelled cells) and uses it to reproduce the three models
of~\cite{severson2019}. The single-feature variance model matches the
published errors to within one cycle, and the six- and nine-feature
elastic-net models agree to within the cross-validation spread, with the
same model ranking (e.g.\ 90 vs.\ 91 cycles RMSE on the primary test set
for the discharge model). The reproduction also exposed a data fault in one
cell that had passed the initial cleaning. Screening 19 physics-informed
candidates showed that most mechanism features are already contained in the
discharge-curve feature, and that several apparent predictors only identify
the test batch. Two features survive: the capacity margin to end of life
divided by the degradation rate, and the early capacity rise caused by the
anode overhang. With the paper's model, these two features beat the
nine-feature model on both test sets (114 vs.\ 125 and 170 vs.\ 220 cycles
RMSE) and equal the six-feature discharge model on the secondary set; they do
not clearly beat it on the primary set, mainly because they under-predict the
longest-lived cell.
A comparison with four other model types shows that,
with 41 training cells, none beats the paper's regularised linear model on a
new batch: a Gaussian process, support vector regression and additive splines
are better within the training batches but not on the secondary batch, and
partial least squares on the whole discharge curve fails there because it
relies on narrow, procedure-dependent features. About 40 cycles suffice for a
rough prediction, 100 cycles are clearly better, and more cycles bring no
gain. On a second unseen batch, not used in the original study, the standard
cells are predicted to within 5--10\,\% by several models, while cells run
with a different procedure are over-predicted by 17--50\,\% by every model;
the model that transfers best differs between the two new batches. SHAP
analysis shows that the models lean on the features that lie outside the
training range, and that a Gaussian process gains over the linear model on the
unseen batch because it learned a saturating effect of the overhang rise.
\end{abstract}

\begin{IEEEkeywords}
lithium-ion battery, cycle life prediction, fast charging, degradation,
machine learning, feature engineering, LFP
\end{IEEEkeywords}
